\section{Twisted bimodules and a lift to the two cones}
\label{sec:lift}

We now construct two stable bimodule maps whose evaluations generate the
one-dimensional degree-three group of the cone calculation.  Their sources
will have distinct twisting actions.  The essential point is a chain-level
lift: we give the homotopies in every degree and identify its projection to
the top cell.

For a nonzero scalar $H\in\kk$, let $h=h_H$ be the automorphism of $T$
that fixes $C$ and multiplies $DC$ by $H$.  Write $T_h$ for the bimodule
with its usual left action and right action $x\cdot a=xh(a)$.  Put
\[
 U=T_h,\qquad Q=P_h,\qquad
 U_H=A_{h\otimes h},\qquad Q^{\mathrm{tot}}=Q\otimes_\kk Q.
\]
Thus $Q$ is the right-twisted bar resolution of $U$, with augmentation
$\epsilon:Q_0\to U$.  A formula on its bar generators extends by
\begin{equation}
 \psi\bigl(a_0[a_1|\cdots|a_i]a_{i+1}\bigr)
 =a_0\psi([a_1|\cdots|a_i])h^{-1}(a_{i+1}).
 \label{lift:extension}
\end{equation}
Here the displayed terminal coefficient is its underlying element of $T$.
This convention makes $\psi$ right linear from the twisted source to an
untwisted target.

\begin{proposition}
\label{lift:main}
For each $H\in\kk^\times$, there is a stable bimodule map
\[
 f_H:U_H\longrightarrow\mathcal C[3]
\]
whose composite with the top-cell projection
$\mathcal C[3]\to A[-1]$ is represented by the bimodule map
$\beta_H:U_H\to\Omega A\subset P_0\otimes P_0$ determined by
\[
 \beta_H(1)=\left(\sum_w h(w)\otimes_I w^*\right)
 \otimes\left(\sum_w h(w)\otimes_I w^*\right),
\]
where each sum is over the letter basis of $T$.
Under the canonical identification $U_H\otimes_A X\simeq X$, its
projection after evaluation sends
\[
 1\longmapsto H^2(f^*\otimes f^*)
 \quad\text{in}\quad
 \Omega X=\rad(A)(f\otimes f)\subset A(f\otimes f).
\]
This evaluated map is nonzero modulo maps factoring through projectives.
In particular, $f_HX$ is nonzero in
$\stHom{A}{X}{\mathcal C X[3]}$.
\end{proposition}

\subsection{A twisted trace-dual tensor}

For a basis letter $w$ of $T$, write $w^*$ for its trace-dual letter.
In $P_0=T\otimes_I T$, set
\begin{equation}
 \xi_H=\sum_{w\text{ in the basis of }T}h(w)\otimes_I w^*.
 \label{lift:casimir}
\end{equation}
The trace-dual coevaluation identity is standard; see \cite[Proposition~3.3]{Broue2009}.
The associative symmetric trace gives the ordinary coevaluation identity
$a\sum_w w\otimes w^*=\sum_w w\otimes w^*a$.
Applying $h$ to the first factor and projecting to $\otimes_I$ gives
\[
 h(a)\xi_H=\xi_Ha.
\]
Consequently $b:U\to P_0$, defined by $b(x)=x\xi_H$, is a bimodule map.

We need the stronger fact that $b$ takes values in $\ker\pi$, where
$\pi:P_0\to T$ is multiplication.  For a lower-case basis letter $w$
of $C$, the product $ww^*$ is the dual of its left vertex idempotent,
and $w^*w$ is the dual of its right vertex idempotent.  Indeed
$(ww^*)(c)=w^*(cw)$; the positive grading makes this zero for
$c\in\rad C$, and the idempotent values determine the assertion.
The row dimensions of $C$ are $6,4$, and its column dimensions are also
$6,4$.  Separating the lower-case and upper-case letters therefore gives
\[
 \pi\xi_H=\sum_{w\text{ lower-case}}ww^*
       +H\sum_{w\text{ lower-case}}w^*w=0
\]
in characteristic two.  The same calculation on each left vertex
$r\in\{e,f\}$ gives
\begin{equation}
 \sum_{\substack{w\in\mathfrak r\text{ a basis letter}\\
                         \mathrm{left}(w)=r}}h(w)w^*=r^*.
 \label{lift:endpoint}
\end{equation}
Here the full sum is zero and the omitted idempotent term is $rr^*=r^*$.

Taking the external square, define
\begin{equation}
 \beta_H=b\otimes_\kk b:
 U_H\longrightarrow\ker(P^{\mathrm{tot}}_0\to A)=\Omega A,
 \qquad P^{\mathrm{tot}}=P\otimes_\kk P.
 \label{lift:beta}
\end{equation}
Its image lies in the indicated kernel because $\pi b=0$.
The remaining task is to lift this map through the top-cell projection.

\subsection{Two homotopies on the bar resolution}

Let $D_b:Q\to P$ be $b\epsilon$ in degree zero and zero in every
positive degree.  It is a chain map.  A homotopy for $D_b$ will supply
a comparison representing $\beta_H$ in the top cell.  Define maps $B_i:Q_i\to P_{i+1}$
for every $i\geq0$ by
\begin{equation}
 B_i([a_1|\cdots|a_i])
 =\sum_{w\in\mathfrak r\text{ a basis letter}}
     [a_1|\cdots|a_i|h(w)]w^*,
 \label{lift:Bformula}
\end{equation}
retaining only composable terms and extending by
\eqref{lift:extension}.

\begin{lemma}
\label{lift:B}
On the nonnegative resolutions, with $B_{-1}=0$, one has
\[
 dB+Bd=D_b.
\]
\end{lemma}
\begin{proof}
For $a\in\mathfrak r$ with right vertex $r$, project the first factor of
$a\xi_H=\xi_Hh^{-1}(a)$ to $\mathfrak r$ along $I$.  Removing the idempotent
terms leaves
\begin{equation}
 \sum_{w\in\mathfrak r}ah(w)\otimes w^*
 +\sum_{w\in\mathfrak r}h(w)\otimes w^*h^{-1}(a)
 =a\otimes r^*,
 \label{lift:projected-casimir}
\end{equation}
with the appropriate vertex restrictions in each sum.
On the left side of the original identity, the omitted term is
$a\otimes r^*$; on its right side the omitted first factors lie in $I$
and disappear under the projection.

For $i>0$, expand $dB_i+B_{i-1}d$ on a bar generator.
All multiplications strictly before the old final bar cancel in pairs.
The two terms adjacent to the old final letter have sum
$[a_1|\cdots|a_i]r^*$ by
\eqref{lift:projected-casimir}.  The last multiplication in $dB_i$ is
exactly the same expression by \eqref{lift:endpoint}, so these terms
also cancel.

At degree zero, on the generator $r[]r$, the differential gives
\[
 dB_0(r[]r)=
 \sum_{\substack{w\in\mathfrak r\\\mathrm{left}(w)=r}}h(w)[]w^*
 +r[]\sum_{\substack{w\in\mathfrak r\\\mathrm{left}(w)=r}}h(w)w^*.
\]
The second term is $r[]r^*$ by \eqref{lift:endpoint}.  It restores the
omitted idempotent term of the full tensor \eqref{lift:casimir}; the
result is $r\xi_H=b\epsilon(r[]r)$.  Bimodule linearity proves the
identity on every degree-zero element.
\end{proof}

We now use $B$ to represent the prescribed top projection.
In $K[3]$, the cell indexed by $\{1,2\}$ is
$P^{\mathrm{tot}}[-1]$.  Put $V=B\otimes D_b$, which raises the
underlying total degree by one, as this shift requires.  Lemma~\ref{lift:B}
gives
\[
 dV+Vd=(b\otimes b)\epsilon^{\mathrm{tot}},
\]
where the right side is supported in degree zero.
Thus $dV_0=\beta_H\epsilon^{\mathrm{tot}}_0$, and at positive indices
$V$ satisfies the comparison equation.  The complex with terms
$P^{\mathrm{tot}}_{i+1}$ and augmentation to $\Omega A$ resolves
$\Omega A$.  Consequently $V$ is a comparison lifting $\beta_H$.

To extend $V$ from the top cell to $K[3]$, we must account for the two
attaching maps.  Deleting slot 1 contributes $pB\otimes D_b$ to the
cell indexed by $\{2\}$.
The suffix comparison $p:P\to P[3]$ from the finite cochain data
lowers the underlying bar degree by three, so $pD_b=0$.
The contribution from deleting slot 2 is therefore zero.
It remains to cancel $pB\otimes D_b$ by a component in cell $\{2\}$.
Lemma~\ref{lift:B} shows that $pB:Q\to P[2]$ is a chain map.
The additional boundary identity~\eqref{coc:boundary} supplies its
nullhomotopy; cocycle closure by itself would not suffice.

\begin{lemma}
\label{lift:G}
There are bimodule maps $G_i:Q_i\to P_{i-1}$ for all $i\geq0$ such that
\[
 pB=dG+Gd.
\]
\end{lemma}
\begin{proof}
Let $z':Q_1\to T$ be the cochain given by $z'([a])=z_H(a)$ from
Lemma~\ref{coc:data}, with extension \eqref{lift:extension}.
Its boundary on $[a|b]$ is
\[
 az'(b)+z'(ab)+z'(a)h^{-1}(b).
\]
By identity~\eqref{coc:boundary}, this equals
$\sum_{w\in\mathfrak r}p(a,b,h(w))w^*$, which is the augmentation of
$(pB)_2$.  Hence
\begin{equation}
 \pi(pB)_2=z'd:Q_2\longrightarrow T.
 \label{lift:boundary}
\end{equation}

Set $G_0=0$ and choose $G_1:Q_1\to P_0$ with $\pi G_1=z'$.
The choice is possible because $Q_1$ is projective and $\pi$ is
surjective.  By \eqref{lift:boundary}, the map
$(pB)_2+G_1d$ has image in $\ker\pi=\im(d:P_1\to P_0)$.
Projectivity of $Q_2$ therefore supplies $G_2$ with
$dG_2=(pB)_2+G_1d$.

Suppose $i\geq3$ and the homotopy equation has been established below
$i$.  Since $pB$ is a chain map,
\[
 d\bigl((pB)_i+G_{i-1}d\bigr)
 =(pB)_{i-1}d+
       \bigl((pB)_{i-1}+G_{i-2}d\bigr)d=0.
\]
Exactness of $P$ in positive degree and projectivity of $Q_i$ allow
this map to be lifted through $d:P_{i-1}\to P_{i-2}$, giving $G_i$.
This proves the equation for every $i\geq2$.  At $i=0,1$ both sides
are zero, since $pB$ has negative target index and $G_0=0$.
\end{proof}

\subsection{The tail map and its evaluation}

The homotopy $G$ supplies the required correction in the singleton cell
$\{2\}$, whose shift in $K[3]$ is $1$.  Together with the top component
$V$, it defines a graded map $\Phi:Q^{\mathrm{tot}}\to K[3]$:
\begin{equation}
 \Phi_{\{1,2\}}=B\otimes D_b,\qquad
 \Phi_{\{2\}}=G\otimes D_b,\qquad
 \Phi_{\{1\}}=\Phi_{\varnothing}=0.
 \label{lift:tail}
\end{equation}
The second component lowers underlying degree by one, as the singleton
shift requires; the first is the comparison $V$ constructed above.

In the top cell, the chain defect $d\Phi+\Phi d$ is
$D_b\otimes D_b$.  In cell $\{2\}$, deletion of slot 1 gives
$pB\otimes D_b$, canceled by
$(dG+Gd)\otimes D_b$.  Deleting slot 2 from either nonzero component
gives zero because $pD_b=0$.  Thus the only defect is
$D_b\otimes D_b$, supported in total source degree zero.
The map therefore commutes with the differential in every positive
total degree.  Passing to high cokernels, at a common index at least $4$ as in
Lemma~\ref{cone:finite}, gives $f_H:U_H\to\mathcal C[3]$.

The top component is still $V$, so the resulting stable projection is
$\beta_H$, as identified above.

\begin{lemma}
\label{lift:evaluation}
Under $U_H\otimes_A X\simeq X$, the evaluated map $\beta_H X$ sends
$1$ to $H^2(f^*\otimes f^*)$ in $\Omega X\subset A(f\otimes f)$.
It does not factor through a projective $A$-module.
\end{lemma}
\begin{proof}
In one slot, tensoring \eqref{lift:casimir} with the simple leaves only
the term whose right coefficient is $f$.  It is the term $w=f^*$, so its
value is $Hf^*$.  Taking the external square gives the stated value.
The cover sequence remains exact on evaluation because it splits on
the right.  Its evaluated degree-zero term is $A(f\otimes f)$, the
projective cover of $X$, with kernel $\rad(A)(f\otimes f)$.

Suppose a map $X\to\rad(A)(f\otimes f)$ factors through a projective.
It then factors through a finite free module.  Every component of a
map $X\to A$ has image in the left socle, since $\rad(A)X=0$.
In a symmetric algebra the left socle also annihilates the radical on
the right: if $\rad(A)u=0$, then, for $j\in\rad(A)$ and $a\in A$,
\[
 \operatorname{tr}(auj)=\operatorname{tr}(jau)=0,
\]
so nondegeneracy of the trace gives $uj=0$.
A map $A\to\rad(A)(f\otimes f)$ is right multiplication by its value
at $1$, which lies in $\rad(A)$, and hence kills the left socle.
Every such factorization is therefore zero.  Our displayed map has
the nonzero value $H^2(f^*\otimes f^*)$, proving the claim.
\end{proof}

By Proposition~\ref{cone:profile}, the top-cell projection is an
isomorphism on the degree-three group of the cone calculation.  Lemma~\ref{lift:evaluation} therefore proves
that $f_HX$ is nonzero, completing Proposition~\ref{lift:main}.
We have obtained the two required maps by taking $H=H_1,H_2$;
the next section combines them into a finite bimodule.
